\section{\textbf{Conclusion}}

We studied world modeling in \textbf{populous, crowded, and chaotic Global
South cities}, where heterogeneity, occlusion, and mixed-traffic negotiation
challenge conventional JEPA predictors. We introduced \textbf{DENSEWORLD}:
\textbf{276 hours} of drive-through, walk-through, and aerial video collected
across \textbf{22 cities}. Building on this benchmark, we proposed
\textbf{\emph{FactorJEPA}}, which makes \textbf{\emph{world structure a
first-class predictive primitive}} through \textbf{\emph{predictor surgery}}: a
three-stage curriculum that adapts the predictor over automatically separated
\textbf{\emph{layout}}, \textbf{\emph{agent}}, and \textbf{\emph{interaction}}
views. Across 2B and 1B backbones, FactorJEPA delivers \textbf{\emph{(i) more
accurate future-latent prediction}}, \textbf{\emph{(ii) stronger causal
future-block structure}}, and \textbf{\emph{(iii) greater robustness under
reduced visual evidence}}, with method rankings reproducing across scale at
$\rho=0.895$ to $0.978$. The results also expose a trade-off between structured
prediction and maximally linear motion readout at the 10k screening scale,
which reverses at full scale (86,831 factor-annotated videos). Together, these
contributions move JEPA world
models from \textbf{\emph{monolithic predictors}} toward
\textbf{\emph{structured, factor-aware models}} of complex real-world dynamics.
